% !TeX program = xelatex
\documentclass[12pt,a4paper,oneside]{report}
% !TeX root = ../main.tex
\XeTeXgenerateactualtext=1
\usepackage[margin=23mm,top=24mm,bottom=24mm,headheight=27pt]{geometry}
\usepackage{amsmath,amssymb,mathtools,bm}
\usepackage{booktabs,array,longtable,tabularx}
\usepackage[table]{xcolor}
\usepackage{tikz}
\usetikzlibrary{arrows.meta,calc,patterns}
\usepackage[most]{tcolorbox}
\usepackage{titlesec,needspace}
\usepackage{hyperref}
\usepackage{bookmark}
\usepackage{xepersian}

% !TeX root = ../main.tex
\settextfont[Scale=1.06,BoldFont=Amiri-Bold.ttf]{Amiri-Regular.ttf}
\setlatintextfont[BoldFont=texgyretermes-bold.otf,ItalicFont=texgyretermes-italic.otf,BoldItalicFont=texgyretermes-bolditalic.otf]{texgyretermes-regular.otf}

% !TeX root = ../main.tex
\definecolor{navy}{HTML}{163747}
\definecolor{teal}{HTML}{127C80}
\definecolor{gold}{HTML}{B18542}
\definecolor{paper}{HTML}{F3F7F7}
\definecolor{muted}{HTML}{5E7178}
\definecolor{alert}{HTML}{A45538}

% مشخصات جزوه؛ عنوان و نام‌ها فقط در این فایل تغییر می‌کنند.
\newcommand{\BookTitle}{داده‌کاوی}
\newcommand{\BookSubtitle}{از دادهٔ خام تا الگو و مدل قابل اعتماد}
\newcommand{\CourseInstructor}{دکتر شیرین شعاعی}
\newcommand{\BookCompilers}{احسان شهبازی}

\hypersetup{unicode=true,colorlinks=true,linkcolor=navy,urlcolor=teal,
 pdftitle={داده‌کاوی ـ کتاب آموزشی و تعاملی},pdfauthor={احسان شهبازی},pdfsubject={پیش‌پردازش، طبقه‌بندی، رگرسیون، خوشه‌بندی و تحلیل وابستگی}}

% !TeX root = ../main.tex
\setlength{\parindent}{0pt}
\setlength{\parskip}{5pt plus 1pt}
\linespread{1.12}
\raggedbottom
\setcounter{tocdepth}{1}
\setcounter{secnumdepth}{2}
\allowdisplaybreaks[2]
\emergencystretch=2em
% Page styles use explicit boxes to keep bidi direction groups balanced.
\makeatletter
\def\ps@lecture{%
  \def\@oddhead{\vbox{\hbox to\textwidth{\small
    \rl{\BookTitle}\hfil\rl{\leftmark}}%
    \kern4pt\hrule height .35pt}}%
  \let\@evenhead\@oddhead
  \def\@oddfoot{\hfil\thepage\hfil}%
  \let\@evenfoot\@oddfoot}
\makeatother
\pagestyle{lecture}
\renewcommand{\chaptermark}[1]{\markboth{#1}{}}
\titleformat{\chapter}[display]{\color{navy}\bfseries\Huge}{\large فصل \thechapter}{8pt}{}[\vspace{4pt}\color{gold}\titlerule]
\titlespacing*{\chapter}{0pt}{0pt}{18pt}
\titleformat{\section}[block]{\Large\bfseries\color{navy}}{\thesection}{.6em}{}
\titleformat{\subsection}[block]{\large\bfseries\color{teal}}{\thesubsection}{.6em}{}

% !TeX root = ../main.tex
\tcbset{enhanced,breakable,boxrule=.45pt,arc=1.5mm,
 left=3mm,right=3mm,top=2mm,bottom=2mm,before skip=9pt,after skip=9pt,
 fonttitle=\bfseries,coltitle=white,colback=paper,colframe=teal}
\newtcolorbox{definition}[1]{breakable=false,title={تعریف | #1},colframe=navy,colback=navy!3!white}
\newtcolorbox{result}[1]{breakable=false,title={نتیجهٔ کلیدی | #1},colframe=teal,colback=teal!4!white}
\newtcolorbox{important}[1]{breakable=false,title={دقت کنید | #1},colframe=gold,colback=gold!6!white}
\newcounter{ex}[chapter]
\newenvironment{example}[1]{\refstepcounter{ex}\begin{tcolorbox}[breakable=false,title={مثال \thechapter.\arabic{ex} | #1},colframe=muted,colback=white]}{\end{tcolorbox}}

% Structured exercises: the web builder turns answers into closed disclosures.
\newenvironment{selftests}{\begin{enumerate}\setlength{\itemsep}{9pt}}{\end{enumerate}}
\newenvironment{selftest}[1]{\item\textbf{#1}\par}{\par}
\newcommand{\answer}[1]{\par{\color{teal}\bfseries پاسخ و توضیح: }#1\par}

% !TeX root = ../main.tex
\newcommand{\edit}[1]{\par\begingroup\small\color{alert}\textbf{تکمیل یا اصلاح علمی.} #1\par\endgroup}
\newcommand{\source}[1]{\par{\small\color{muted} منبع مطالعه: #1}\par}
\newcommand{\E}{\mathbb E}
\newcommand{\Var}{\operatorname{Var}}
\newcommand{\Cov}{\operatorname{Cov}}
\newcommand{\MSE}{\operatorname{MSE}}
\newcommand{\Bias}{\operatorname{Bias}}
\newcommand{\SE}{\operatorname{SE}}
\newcommand{\iid}{\overset{\mathrm{iid}}{\sim}}
\newcommand{\N}{\mathcal N}
\newcommand{\R}{\mathbb R}
\newcommand{\Prob}{\mathbb P}
\newcommand{\sol}{\par\textbf{حل.} }
\newcommand{\figcap}[1]{\par\smallskip{\small\color{muted}#1}\par}

\newcommand{\labref}[1]{\par\smallskip{\small\color{teal}\textbf{آزمایشگاه پیشنهادی.} #1}\par}
\newcommand{\checkanswer}[1]{\par\textbf{پاسخ کوتاه.} #1\par}

\providecommand{\tightlist}{\setlength{\itemsep}{0pt}\setlength{\parskip}{0pt}}
\providecommand{\passthrough}[1]{#1}



\hypersetup{pdftitle={طبقه‌بندی؛ درخت، تجمیع و ارزیابی},pdfauthor={دکتر شیرین شعاعی؛ احسان شهبازی}}
\begin{document}
\setcounter{chapter}{2}
\chapter{طبقه‌بندی؛ درخت، تجمیع و ارزیابی}
\begin{center}{\small\color{muted} استاد درس: دکتر شیرین شعاعی\\[2mm] گردآورنده جزوه: احسان شهبازی}\end{center}

\source{اسلاید طبقه‌بندی ۱، PDF ۳ تا ۶۰؛ طبقه‌بندی ۲، PDF ۴۲ تا ۶۱؛ کتاب فصل ۳، صفحات چاپی ۱۳۴ تا ۱۹۳، و بخش‌های ۶٫۱۰ و ۶٫۱۱.}
\section{از مثال‌های برچسب‌دار تا تصمیم برای دادهٔ تازه}

در طبقه‌بندی (Classification) دادهٔ آموزشی شامل زوج‌های $(\mathbf x_i,y_i)$ است؛ $\mathbf x_i$ بردار ویژگی‌ها و $y_i$ برچسب گسسته است. الگوریتم باید تابعی مانند $f(\mathbf x)$ بسازد تا برای یک مشاهدهٔ تازه برچسب پیش‌بینی کند. تفاوت با رگرسیون در نوع خروجی است، نه در عددی بودن ورودی: پیش‌بینی قیمت خانه رگرسیون است و پیش‌بینی «پرخطر/کم‌خطر» بودن وام طبقه‌بندی. طبقه‌بندی دوکلاسه می‌تواند «تقلب/سالم» و چندکلاسه می‌تواند نوع گیاه باشد. نمونه‌های آموزشی هرگز تعریف کامل مفهوم نیستند؛ مدلی که فقط درخت‌های سبز و بزرگ دیده ممکن است گیاه بزرگ گلدانی را هم درخت تشخیص دهد.

گردش کار صحیح از تعریف مسئله و واحد مشاهده شروع می‌شود؛ سپس دادهٔ برچسب‌دار به بخش‌های مناسب جدا می‌شود، پیش‌پردازش فقط از آموزش آموخته می‌شود، مدل آموزش می‌بیند، تنظیمات با اعتبارسنجی انتخاب می‌شوند و آزمون مستقل کیفیت تعمیم را می‌سنجد. هدف صرفاً حفظ داده‌های گذشته نیست. «یادگیری مشتاق» (Eager) پیش از دریافت پرسش مدل صریح می‌سازد؛ درخت تصمیم نمونهٔ روشن آن است. روش تنبل (Lazy) مانند kNN بیشتر نمونه‌ها را نگه می‌دارد و بخش مهم محاسبات را به زمان پیش‌بینی موکول می‌کند. تنبل بودن به معنی بی‌هزینه بودن کل سامانه نیست و مشتاق بودن نیز تفسیرپذیری را تضمین نمی‌کند؛ شبکهٔ عصبی مشتاق است ولی معمولاً توضیح انسانی ساده ندارد.



\section{درخت تصمیم: پرسش‌های ساده، مرزهای مستطیلی}

درخت تصمیم (Decision Tree) زنجیره‌ای از آزمون‌هاست. ریشه نخستین پرسش است، هر گره داخلی آزمون ویژگی، هر یال یک پاسخ و هر برگ تصمیم نهایی است. پیش‌بینی با عبور از ریشه به برگ انجام می‌شود. در مثال اسلاید، ابتدا وضعیت بازپرداخت مالیات پرسیده می‌شود؛ سپس وضعیت تأهل و در نهایت درآمد مشمول مالیات. برای رکورد «\lr{Refund=No} ، متأهل، درآمد۸۰هزار» مسیر به برگ \lr{Cheat=No} می‌رسد و نیازی نیست همهٔ ویژگی‌ها در هر پیش‌بینی بررسی شوند.

اگر هر گره فقط یک ویژگی را با آستانه مقایسه کند، مرزها موازی محورند؛ $x<0.43$ یک خط عمودی و $y<0.33$ خط افقی می‌سازد. ترکیب تعداد زیادی چنین مستطیل‌هایی می‌تواند مرز خمیده را تقریب بزند، اما ممکن است درخت بزرگ شود. این محدودیت مربوط به درخت تک‌ویژگی متداول است؛ همهٔ انواع درخت در جهان الزاماً چنین نیستند (کتاب در پایان۳.۳.۶ درخت مایل را نیز نام می‌برد).

الگوریتم هانت (Hunt) تقسیم بازگشتی حریصانه دارد: اگر همهٔ برچسب‌ها یکی باشند برگ می‌سازیم؛ وگرنه آزمون‌های ممکن را مقایسه می‌کنیم، آزمون دارای بیشترین کاهش ناخالصی را انتخاب و داده‌ها را میان فرزندان تقسیم می‌کنیم؛ همین کار برای هر فرزند تکرار می‌شود. اگر ویژگی‌ها یکسان ولی برچسب‌ها متناقض باشند تقسیم ممکن نیست و برگ با رأی اکثریت برچسب می‌گیرد. اگر یک شاخه هیچ مشاهده‌ای ندارد می‌توان اکثریت والد را به کار برد. این روش بهینهٔ سراسری را تضمین نمی‌کند؛ بهترین پرسش محلی الزاماً بهترین درخت نهایی نیست.

برای ویژگی اسمی با k مقدار، تقسیم چندشاخه یک شاخه برای هر مقدار می‌سازد؛ تقسیم دودویی همهٔ مقادیر را در دو زیرمجموعه قرار می‌دهد و تعداد تقسیم‌های غیرتکراری غیرتهی $2^{k-1}-1$ است. برای ویژگی ترتیبی باید ترتیب حفظ شود؛ \{کوچک،متوسط\}/\{بزرگ\} معتبر است، ولی \{کوچک،بزرگ\}/\{متوسط\} ترتیب را به هم می‌زند. برای ویژگی پیوسته معمولاً دو شاخهٔ $A\le v$ و $A>v$ کافی است؛ مقدار دقیق برابری باید صریح باشد تا هیچ رکوردی گم یا دوبار شمرده نشود.

\section{ناخالصی جینی و محاسبهٔ یک تقسیم}

اگر $p_j(t)$ سهم طبقهٔ $j$ در گره $t$ باشد، ناخالصی جینی چنین تعریف می‌شود:
\[G(t)=1-\sum_{j=1}^c p_j(t)^2.\]
گره خالص جینی صفر دارد؛ بیشینه برای $c$ طبقهٔ هم‌احتمال $1-1/c$ است. در مسئلهٔ دوطبقه‌ای $G=2p(1-p)$ و بیشینه $0.5$ است. برای شمارش‌های $(1,5)$ مقدار $10/36\approx0.2778$ و برای $(2,4)$ مقدار $16/36\approx0.4444$ به دست می‌آید. این معیار با ضریب جینی نابرابری درآمد یک تعریف ندارد.

پس از شکستن، اندازهٔ فرزندان در وزن‌دهی ضروری است:
\[G_{\mathrm{split}}=\sum_{r=1}^m\frac{n_r}{n}G(t_r),\qquad \Delta G=G(t)-G_{\mathrm{split}}.\]
برای والد $(6,6)$ و فرزندان $(5,2)$ و $(1,4)$، جینی والد $0.5$، جینی فرزندان $20/49$ و $8/25$ و جینی وزنی $13/35\approx0.37143$ است. پس سود خلوص $9/70\approx0.12857$ می‌شود. استفاده از میانگین سادهٔ دو فرزند، سهم جمعیت هر شاخه را نادیده می‌گیرد.

کتاب معیار آنتروپی و سود اطلاعاتی را نیز بسط می‌دهد:
\[H(t)=-\sum_j p_j(t)\log_2p_j(t),\qquad 0\log 0=0.\]
کاهش آنتروپی، سود اطلاعاتی و کاهش جینی، سود خلوص است. یک شناسهٔ یکتا می‌تواند همهٔ برگ‌های آموزش را خالص کند، ولی به‌تنهایی تعمیم را بهتر نمی‌کند. نسبت سود، سود اطلاعاتی را نسبت به آنتروپیِ اندازهٔ شاخه‌ها می‌سنجد و یکی از راه‌های کنترل تمایل به تقسیم‌های پرشاخه است.

\section{پیدا کردن آستانهٔ ویژگی پیوسته}

داده‌ها را بر حسب ویژگی مرتب می‌کنیم، شمار کلاس‌های چپ و راست را نگه می‌داریم، و هنگام عبور آستانه فقط شمارش رکورد جابه‌جا‌شده را به‌روز می‌کنیم. مرتب‌سازی $O(n\log n)$ و پیمایش پس از آن $O(n)$ است؛ اسکن مجدد کل داده برای هر آستانه $O(n^2)$ می‌شود. نقاط میانی میان مقادیر متمایز کافی‌اند و مقادیر تکراری باید به صورت گروهی جابه‌جا شوند. در مثال درآمدهای ۶۰،۷۰،۷۵،۸۵،۹۰،۹۵،۱۰۰،۱۲۰،۱۲۵،۲۲۰ با برچسب‌های \lr{No, No, No, Yes, Yes, Yes, No, No, No, No} آستانه۹۷٫۵ جینی۰٫۳ دارد؛ آستانه۸۰ جینی۰٫۳۴۳ می‌دهد. حذف نامزدهای بین برچسب‌های یکسان در حالت استاندارد بدون محدودیت اضافه محاسبات را کم می‌کند؛ با وزن‌ها یا قیود خاص باید احتیاط کرد.

ابهام مهم مثال هانت اسلاید۱۲: گفته شده Refund خالص‌ترین تقسیم اولیه است، اما در جدول نامزدهای پیوستهٔ خود اسلاید۲۸، بهترین آستانهٔ درآمد $G_{\mathrm{split}}=0.300$ دارد و از Refund با۰٫۳۴۳ بهتر است. بنابراین روایت هانت مثالِ روند ساخت است، نه اثبات اینکه آن ریشه در میان تمام آستانه‌ها بهینه است. این دو مثال، هدف و مجموعهٔ نامزدهای یکسانی را اثبات نمی‌کنند.



\section{بیش‌برازش، سادگی و تجمیع مدل‌ها}

دو درخت در مثال اعتبار مالی همهٔ پنج رکورد آموزش را درست پیش‌بینی می‌کنند. اولی می‌پرسد بدهی از۵۰۰۰ بیشتر است یا نه؛ دومی فقط نام John و Alice را حفظ می‌کند. عملکرد آموزش مساوی است ولی شواهدی نداریم که حفظ نام‌ها برای مشتری جدید مفید باشد. اصل تیغ اُکام می‌گوید میان توضیح‌هایی با کیفیت مشابه، توضیح ساده‌تر را ترجیح دهیم؛ این اصل یک راهنمای انتخاب است و تضمین همیشگی برتری هر مدل کوچک نیست.

کم‌برازش یعنی مدل حتی رابطهٔ اصلی را نمی‌گیرد؛ درخت تک‌برگ برای مسئلهٔ پیچیده مثال آن است. با افزایش ظرفیت ابتدا خطای آموزش و اعتبارسنجی هر دو کم می‌شوند؛ پس از نقطه‌ای آموزش باز هم بهتر ولی اعتبارسنجی بدتر می‌شود. بیش‌برازش یعنی یادگیری جزئیات تصادفی نمونه به جای ساختار قابل‌تعمیم. نویز، نمونهٔ کم، ویژگی‌های نامرتبط زیاد و ظرفیت بالا زمینهٔ آن را فراهم می‌کنند. عمق زیاد علامت قطعی نیست؛ عمق را نسبت به مقدار داده و پیچیدگی واقعی مسئله باید سنجید. در نمودار آموزشی می‌توان خطای آزمون را برای تشریح پدیده نشان داد، اما تنظیم واقعی مدل باید با اعتبارسنجی انجام شود.

سه راه موردتأکید اسلاید عبارت‌اند از دادهٔ آموزشی بیشتر، پیش‌هرس و تجمیع. دادهٔ بیشتر وقتی مفید است که نمایندهٔ جمعیت هدف باشد؛ کپی چندبارهٔ همان رکورد اطلاعات تازه نمی‌دهد. پیش‌هرس با max\_depth ، حداقل اندازهٔ برگ یا حداقل کاهش ناخالصی رشد را محدود می‌کند. توقف بسیار زود نیز خطر کم‌برازش دارد: در XOR ممکن است تقسیم اول سود اندک بدهد ولی تقسیم بعدی لازم باشد. کتاب پس‌هرس را نیز اضافه می‌کند: نخست درخت بزرگ ساخته و سپس زیر‌درخت‌های بی‌فایده با برگ یا شاخهٔ ساده‌تر جایگزین می‌شوند؛ این بسط از کتاب است نه موضوع مفصل اسلاید.

تجمیع \lr{Ensemble} رأی چند مدل پایه را ترکیب می‌کند. اگر همه یک اشتباه کنند، رأی‌گیری کمکی ندارد. برای $M$ رأی‌دهندهٔ مستقل، با $M$ فرد و احتمال خطای یکسان $q$، احتمال خطای اکثریت برابر است با
\[\sum_{i=(M+1)/2}^{M}\binom Mi q^i(1-q)^{M-i}.\]
این مدل ساده نقش تنوع خطاها را روشن می‌کند؛ در کاربرد، استقلال کامل معمول نیست و این فرمول تضمین عملکرد جنگل تصادفی نیست.

در Bagging از دادهٔ آموزش نمونه‌گیری با جایگذاری می‌کنیم و برای هر نمونه یک درخت جدا می‌سازیم. یک نمونهٔ بوت‌استرپ n تایی فقط حدود۶۳٫۲٪ رکورد متمایز دارد، زیرا احتمال حداقل یک بار انتخاب رکورد$1-(1-1/n)^n$ است. حدود۳۶٫۸٪ برای همان درخت خارج از کیسه می‌مانند. جنگل تصادفی علاوه بر بوت‌استرپ، در هر گره تنها یک زیرمجموعهٔ تصادفی ویژگی را نامزد تقسیم می‌کند؛ این کار غلبهٔ همیشگی یک ویژگی قوی را می‌شکند و همبستگی درخت‌ها را کاهش می‌دهد. «استقلال درخت‌ها» در اسلاید توصیف ایده‌آل است؛ جنگل استقلال دقیق آماری تضمین نمی‌کند. p خیلی کوچک درخت را ضعیف می‌کند، p بزرگ جنگل را شبیه Bagging می‌کند. تعداد درخت‌ها و اندازهٔ زیرمجموعهٔ ویژگی‌ها ابرپارامترند. درخت تنها معمولاً تفسیرپذیرتر است؛ جنگل پیچیده‌تر و اغلب قوی‌تر، ولی برتری همیشگی ندارد.

بوستینگ مدل را افزایشی می‌سازد؛ مدل تازه برای کاهش زیان مجموعهٔ موجود آموزش می‌بیند:
\[F_t(\mathbf x)=F_{t-1}(\mathbf x)+\eta f_t(\mathbf x).\]
در معرفی XGBoost در اسلاید، تابع هدف شامل زیان و جریمهٔ پیچیدگی درخت است:
\[\operatorname{Obj}=L+\Omega,\qquad \Omega(f)=\gamma T+\frac\lambda2\sum_{j=1}^{T}w_j^2.\]
$T$ تعداد برگ و $w_j$ امتیاز برگ است. برای سه برگ با امتیازهای $2,0.1,-1$، جریمه $3\gamma+2.505\lambda$ است؛ برای هر دو ضریب برابر یک، مقدار $5.505$ می‌شود. ساختار گسستهٔ درخت مانند وزن‌های پیوستهٔ شبکه با یک گام معمول گرادیان حرکت نمی‌کند؛ گرادیان زیان برای ساخت درخت بعدی استفاده می‌شود. این با ساخت درخت‌های جدا در جنگل تصادفی متفاوت است.



\section{ارزیابی: کدام اشتباه و با چه هزینه‌ای؟}

ابتدا کلاس مثبت را روشن تعریف می‌کنیم؛ «مثبت» به معنی خوب یا مطلوب نیست، در تشخیص تقلب همان تقلب مثبت است. TP مثبت واقعیِ درست شناسایی‌شده، FN مثبت واقعیِ ازدست‌رفته، FP منفی واقعی با هشدار مثبت و TN منفی واقعیِ درست تشخیص‌داده‌شده است. در این جزوه سطر حقیقت و ستون پیش‌بینی است. حرف اول True/False درستی پیش‌بینی و حرف دوم Positive/Negative خروجی مدل را توصیف می‌کند.

با قرارداد سطرِ حقیقت و ستونِ پیش‌بینی، $n=TP+TN+FP+FN$ است. صحت کلی و نرخ خطا چنین‌اند:
\[\operatorname{Accuracy}=\frac{TP+TN}{n},\qquad \operatorname{Error}=1-\operatorname{Accuracy}.\]
در مثال $TP=25,FN=4,FP=6,TN=15$، صحت $40/50=0.8$ است. باید آن را با مبنا سنجید: پیش‌بینی همیشگی طبقهٔ اکثریت در همین داده $29/50=0.58$ صحت دارد. درصد بدون اندازهٔ آزمون، خط مبنا و هزینهٔ خطا، قضاوت کاملی نیست.

در داده‌ای با ۹۹۹۰ منفی و ۱۰ مثبت، پیش‌بینی همیشگی منفی صحت ۹۹٫۹ درصد دارد ولی هیچ مثبت را نمی‌یابد. درستی هشدار و بازیابی دو سؤال متفاوت‌اند:
\[P=\operatorname{Precision}=\frac{TP}{TP+FP},\qquad R=\operatorname{Recall}=\frac{TP}{TP+FN}.\]
اولی سهم هشدارهای صحیح و دومی سهم مثبت‌های واقعیِ کشف‌شده است. در مثال پیشین $P=25/31\approx0.80645$ و $R=25/29\approx0.86207$ است.

میانگین هارمونیک درستی و بازیابی، $F_1$، به کم بودن یکی از دو معیار حساس است:
\[F_1=\frac{2PR}{P+R}=\frac{2TP}{2TP+FP+FN}.\]
در مثال بالا $F_1=50/60\approx0.83333$ است. صورت عمومیِ صحیح چنین است:
\[F_\beta=\frac{(1+\beta^2)PR}{\beta^2P+R}=\frac{(1+\beta^2)TP}{(1+\beta^2)TP+\beta^2FN+FP}.\]
$\beta>1$ بازیابی را مهم‌تر می‌کند. برای $TP=1,FP=0,FN=99$، درستی کامل ولی بازیابی $0.01$ و $F_1=2/101$ است. اگر مخرج یک معیار صفر باشد باید قرارداد گزارش روشن باشد؛ صفر قراردادی و مقدار تعریف‌نشده از نظر مفهومی یکسان نیستند.

هزینهٔ تصمیم از شمار درست و غلط فراتر می‌رود. در اسلاید ۵۸، هزینه‌های $TP,FN,FP,TN$ به‌ترتیب $-1,100,1,0$ هستند. هزینهٔ منفی پاداش است. مدل اول شمارش‌های $(162,38,160,240)$ و مدل دوم $(155,45,5,395)$ را با همین ترتیب دارد:
\begin{align*}
\operatorname{Accuracy}_1&=\frac{402}{600}=0.67,\\
\operatorname{Cost}_1&=-162+3800+160=3798,\\
\operatorname{Accuracy}_2&=\frac{550}{600}\approx0.9167,\\
\operatorname{Cost}_2&=-155+4500+5=4350.
\end{align*}
بنابراین با این هزینه‌ها مدل اول مطلوب‌تر است، با اینکه صحت پایین‌تری دارد. هزینهٔ کل $\sum_{a,p}n_{a,p}c_{a,p}$ است؛ تقسیم آن بر تعداد مشاهده‌ها، هزینهٔ متوسط را می‌دهد. این دو مدل روی تعداد مثبت و منفی یکسان مقایسه شده‌اند.

\section{امتیاز، آستانه و ROC}

بسیاری مدل‌ها به جای برچسب فقط، امتیاز مثبت‌بودن $s(x)$ می‌دهند. قانون $s(x)\ge\tau$ برچسب مثبت می‌سازد؛ با کاهش $\tau$ افراد بیشتری مثبت اعلام می‌شوند و Recall هرگز کاهش نمی‌یابد، ولی Precision الزاماً یکنواخت تغییر نمی‌کند. عبارت اسلاید۵۶ «آستانهٔ بالاتر همیشه Precision بیشتر» یک روند معمول است، نه قضیه برای هر نمونه. آستانه را روی اعتبارسنجی متناسب با هدف یا هزینه انتخاب می‌کنیم، نه با آزمون نهایی.

در نمودار ROC ، نرخ مثبت واقعی در برابر نرخ مثبت کاذب رسم می‌شود:\[\operatorname{TPR}=\frac{TP}{TP+FN},\qquad\operatorname{FPR}=\frac{FP}{FP+TN}.\] از آستانهٔ بالاتر از همهٔ امتیازها در(۰،۰) شروع و امتیازها را نزولی بررسی می‌کنیم. عبور از یک مثبت واقعی گام عمودی $1/n_{\mathrm{positive}}$ و عبور از یک منفی واقعی گام افقی $1/n_{\mathrm{negative}}$ می‌دهد. اسلاید۵۹ «پیش‌بینی درست/نادرست» می‌گوید؛ دقیق‌تر و صحیح‌تر «برچسب واقعی مثبت/منفی هنگام جاروب امتیاز مثبت» است. اگر امتیازها مساوی‌اند باید همزمان پردازش شوند؛ ترتیب دلخواه میان هم‌امتیازها AUC را تحریف می‌کند. نقطهٔ ایده‌آل(۰،۱)، پیش‌بینی همگی مثبت(۱،۱) و خط انتظار تصادفی قطر است.

\begin{example}{محاسبهٔ ROC با چهار مشاهده}
امتیازهای نزولی $0.9,0.8,0.7,0.4$ و برچسب‌های $+,-,+,-$ را در نظر بگیرید. نقاط ROC به‌ترتیب چنین‌اند:
\[(0,0),\ (0,0.5),\ (0.5,0.5),\ (0.5,1),\ (1,1).\]
مساحت ذوزنقه‌ها $0.75$ است؛ از چهار جفت مثبت–منفی، سه جفت ترتیب درست دارند. در آستانهٔ $0.7$، شمارش‌ها $TP=2,FP=1,FN=0,TN=1$ و درستی $2/3$ و بازیابی یک‌اند. AUC کیفیت رتبه‌بندی را خلاصه می‌کند، نه کالیبراسیون یا احتمال صحیح بودن تک‌پیش‌بینی‌ها. برای عدم توازن شدید، منحنی درستی–بازیابی نیز اهمیت دارد.
\end{example}



\section{ارزیابی معتبر، انتخاب ابرپارامتر و نشت اطلاعات}

معیار ارزیابی «چه عددی بسنجیم» را تعیین می‌کند، روش ارزیابی «این عدد را از کدام داده و چه آزمایشی به دست آوریم» را. Holdout یک بخش مستقل را برای آزمون کنار می‌گذارد؛ نسبت۸۰/۲۰ یا۲/۳--۱/۳ قانون همگانی نیست. آموزش کوچک مدل را ضعیف می‌کند و آزمون کوچک برآورد را پرنوسان. در نامتوازنی، نمونه‌گیری طبقه‌بندی‌شده (Stratified) نسبت کلاس‌ها را تقریباً حفظ می‌کند، نه اینکه تعداد مثبت و منفی را مساوی کند. برای داده‌های زمانی یا چند رکورد از یک شخص، جداسازی زمانی/گروهی بر تصادفی‌سازی ساده مقدم است؛ این قید تکمیلی برای جلوگیری از نشت است.

در Repeated Holdout هر بار تقسیم تصادفی تازه و امتیازها میانگین گرفته می‌شوند؛ آزمون‌ها ممکن است همپوشانی داشته باشند و برخی نمونه‌ها هرگز آزمون نشوند. در k-fold داده به k بخش تقسیم می‌شود و هر بار یکی آزمون و بقیه آموزش است؛ هر مشاهده یک بار آزمون و $k-1$ بار آموزش می‌شود. با اندازه‌های برابر، میانگین خطای fold ها مناسب است؛ برای اندازه‌های متفاوت، خطای شمارشی کل تقسیم بر تعداد کل دقیق‌تر است. اگر از F1 استفاده می‌کنیم، میانگین F1fold ها و F1 حاصل از تجمیع پیش‌بینی‌های out-of-fold لزوماً یکی نیستند؛ روش گزارش باید صریح باشد.

در Leave-One-Out ، $k=n$ است و هر بار یک مشاهده آزمون می‌شود. دادهٔ آموزش زیاد است ولی هزینهٔ آموزش n مدل و گاه واریانس بالا دارد. جملهٔ اسلاید۴۸ «برآورد بسیار دقیق» تضمین همگانی نیست؛ خود کتاب محدودیت‌ها و سناریوهای گمراه‌کننده را گوشزد می‌کند. برای $n=10^6$ و هر آموزش۱دقیقه، حدود۶۹۴روز زمان سریال لازم می‌شود. مقادیر۵یا۱۰ fold مصالحه‌های رایج معرفی‌شده‌اند، نه نتیجهٔ بهینهٔ قطعی برای هر مسئله.

پارامتر از داده یاد گرفته می‌شود: وزن شبکه، میانگین بیز و محل تقسیم درخت. ابرپارامتر قواعد آموزش یا معماری را تعیین می‌کند: عمق، C ،$\gamma$،$\eta$،تعداد لایه. انتخاب ابرپارامتر نیز نوعی یادگیری از دادهٔ اعتبارسنجی است و نباید از آزمون نهایی انجام شود. جست‌وجوی شبکه‌ای تمام حاصل‌ضرب دکارتی مقادیر را می‌آزماید؛ جست‌وجوی تصادفی از توزیع‌های مشخص نمونه می‌گیرد و وقتی فقط چندبعد حساس‌اند اغلب پوشش مؤثرتری می‌دهد. برتری همیشگی تصادفی بر شبکه‌ای ادعا نمی‌شود. بهینه‌سازی تکاملی ترکیب‌های مناسب را نگه و تغییر می‌دهد؛ بهینه‌سازی بیزی یک مدل جانشین از امتیاز برحسب ابرپارامتر می‌سازد و با درنظرگرفتن عدم‌قطعیت نامزد بعدی را پیشنهاد می‌کند. واژهٔ «بیزی» اینجا ارتباطی با انتخاب طبقه‌بند Naïve Bayes ندارد.

روش سه‌بخشی:۱) آزمون را قفل کنید؛۲)آموزش را به train/validation تقسیم کنید؛۳)هر تنظیم را فقط با train آموزش و با validation بسنجید؛۴)بهترین تنظیم را انتخاب کنید؛۵)با آن تنظیم روی کل دادهٔ غیرآزمون دوباره آموزش دهید؛۶)یک بار آزمون نهایی انجام دهید. مثلاً امتیازهای اعتبارسنجی F1 برای سه تنظیم۰٫۶،۰٫۸،۰٫۷‌اند؛ تنظیم دوم انتخاب می‌شود، اما۰٫۸ تخمین نهایی تعمیم نیست. نمونهٔ اسلاید پس از آموزش مجدد F1آزمون۰٫۶۸ دارد؛ کاهش آن تناقض نیست.

اگر اعتبارسنجی ساده ناپایدار باشد، CV داخلی برای انتخاب تنظیم به کار می‌رود. اگر کل داده باید هم در انتخاب و هم در برآورد تعمیم به‌صورت منظم استفاده شود، CV تودرتو (Nested) لازم است: در هر fold بیرونی، آزمون بیرونی کاملاً کنار می‌ماند؛ فقط دادهٔ آموزش بیرونی به CV داخلی می‌رود؛ تنظیم برنده روی کل آموزش بیرونی آموزش مجدد می‌بیند و روی آزمون بیرونی سنجیده می‌شود. میانگین امتیازهای بیرونی کیفیت «رویهٔ انتخاب و آموزش» را برآورد می‌کند، نه الزاماً یک مدل خاص با تنظیم ثابت. پس از گزارش ارزیابی می‌توان برای مدل قابل‌استفاده تنظیم را روی کل داده انتخاب و مدل نهایی را آموزش داد؛ این مرحله آزمون تازه‌ای تولید نمی‌کند.

برای ۱۰۰ تنظیم و ۵ بخش، تعداد آموزش‌ها در سه شیوه چنین است:
\begin{align*}
\text{Holdout}:&\quad100+1=101,\\
\text{CV}:&\quad5\times100+1=501,\\
\text{Nested CV}:&\quad5(5\times100+1)=2505.
\end{align*}
حالت آخر ۵ بخش بیرونی و ۵ بخش داخلی دارد. این شمارش، برازش مجدد هر بخش بیرونی را دارد و آموزش مدل نهایی روی تمام داده را شامل نمی‌شود. تعداد بیشتر بدون طراحی مستقل می‌تواند فقط بیش‌برازشِ انتخاب را بیشتر کند.

انتخاب ویژگی Forward از بهترین تک‌ویژگی شروع و ویژگی تازه اضافه می‌کند؛ Backward از همه شروع و حذف می‌کند. کل این جست‌وجو، تخمین مقیاس، imputation ، PCA وهرگونه توازن‌سازی باید درون آموزش هر fold باشد. حتی درخت، SVM یا ANN که می‌توانند وزن کم به ویژگی بدهند، الزاماً انتخاب زیرمجموعهٔ دقیق و خودکار انجام نمی‌دهند؛ L2SVM غالباً وزن‌های دقیقاً صفر نمی‌سازد. کاهش اثر ویژگی الزاماً انتخاب یک زیرمجموعهٔ دقیق نیست. در دادهٔ نامتوازن بازنمونه‌گیری فقط آموزش است؛ نسبت آزمون باید مسئلهٔ واقعی را بازنمایی کند. لزوم Nested به وابستگی انتخاب و ارزیابی مربوط است؛ «تکرارپذیری فقط با یک split ممکن است» درست نیست و seed/نسخه/داده را برای CV هم می‌توان ثبت کرد.



\subsection*{سه سطح داده در اعتبارسنجی تودرتو}
در یک دور بیرونی، آزمون بیرونی را کنار می‌گذاریم. برای هر تنظیم نامزد، چند دور داخلی روی دادهٔ باقی‌مانده اجرا می‌شود؛ هر دور داخلی آموزش، جایگذاری، مقیاس‌بندی و مدل خاص خود را دارد. پس از انتخاب تنظیم، مدل از ابتدا با کل دادهٔ غیرآزمون همان دور بیرونی برازش و فقط یک بار روی آزمون بیرونی سنجیده می‌شود. آزمون بیرونی نباید برای توقف زودهنگام یا انتخاب ویژگی استفاده شود. با ۵ دور بیرونی، ۴ دور داخلی و ۱۰ تنظیم:
\[N_{\mathrm{fit}}=5(4\times10+1)=205.\]
این عدد شامل بازبرازش هر دور بیرونی است. میانگین امتیازهای بیرونی، کیفیت رویهٔ انتخاب را توصیف می‌کند؛ اختلاف جزئی امتیازها بدون توجه به تغییرپذیری داده، دلیل قطعی برتری نیست.

\section{خودآزمایی و تمرین}
\begin{selftests}
\begin{selftest}{جینیِ یک تقسیم را محاسبه کنید}
والد با شمارش $(5,5)$ به دو فرزند $(4,1)$ و $(1,4)$ تقسیم می‌شود. جینی وزنی و سود چقدر است؟
\answer{جینی والد $0.5$ و جینی هر فرزند $1-(4/5)^2-(1/5)^2=0.32$ است. اندازه‌ها برابرند، پس جینی وزنی $0.32$ و سود $0.18$ است. با اندازه‌های نابرابر، میانگین ساده جای وزن‌دهی را نمی‌گیرد.}
\end{selftest}
\begin{selftest}{صحت زیاد را نقد کنید}
در ۱۰۰۰ مشاهده با ۱۰ مثبت، مدلی همیشه منفی می‌گوید. صحت و بازیابی چیست؟
\answer{صحت $990/1000=0.99$ ولی بازیابی صفر است. چون هیچ هشدار مثبتی صادر نشده، درستی هشدار مخرج صفر دارد و تعریف‌نشده است؛ نباید آن را با درستی کامل اشتباه گرفت.}
\end{selftest}
\begin{selftest}{آستانه و AUC را جدا کنید}
با ثابت بودن امتیازها، کاهش آستانه چه اثری بر بازیابی و AUC دارد؟
\answer{بازیابی کاهش نمی‌یابد، زیرا هیچ مثبتِ قبلاً کشف‌شده‌ای از مجموعهٔ هشدارها حذف نمی‌شود. AUC از تمام آستانه‌ها و ترتیب امتیازها حاصل می‌شود و با حرکت نقطهٔ جاری تغییر نمی‌کند. درستی هشدار الزاماً یکنواخت نیست.}
\end{selftest}
\begin{selftest}{تجمیع چه زمانی کمکی نمی‌کند؟}
اگر همهٔ درخت‌ها روی دقیقاً همان مشاهده‌ها خطا کنند، رأی اکثریت چه کمکی می‌کند؟
\answer{برای آن خطاهای مشترک کمکی نمی‌کند. تنوع خطاها مهم است؛ جنگل تصادفی با بوت‌استرپ و نامزدی تصادفی ویژگی‌ها همبستگی را کاهش می‌دهد، اما استقلال کامل یا برتری همیشگی را تضمین نمی‌کند.}
\end{selftest}
\begin{selftest}{داده‌های هر دور را مشخص کنید}
در طرح ۵ بیرونی، ۴ داخلی و ۱۰ تنظیم، چند برازش داریم؟ آیا آزمون بیرونی برای توقف زودهنگام مجاز است؟
\answer{با بازبرازش هر دور بیرونی، ۲۰۵ برازش داریم. استفاده از آزمون بیرونی برای توقف یا تنظیم، استقلال آن را از بین می‌برد؛ این تصمیم‌ها باید در دادهٔ آموزش بیرونی انجام شوند.}
\end{selftest}
\begin{selftest}{خطای آموزش کافی است؟}
خطای آموزش با افزایش عمق از ۱۰ درصد به صفر می‌رسد، ولی اعتبارسنجی از ۱۵ به ۲۵ درصد می‌رود. چه چیزی را تغییر می‌دهید؟
\answer{این الگو با بیش‌برازش سازگار است؛ عمق یا حداقل اندازهٔ برگ را با اعتبارسنجی انتخاب می‌کنیم و نشت و نمایندگی تقسیم را بررسی می‌کنیم. آموزش را برای بهتر شدن امتیاز آزمون نهایی تکرار نمی‌کنیم.}
\end{selftest}
\begin{selftest}{واحدِ درآمد چه اثری بر درخت دارد؟}
چرا تبدیل ریال به میلیون ریال معمولاً تقسیم‌های درخت تک‌ویژگی را عوض نمی‌کند؟
\answer{تبدیل مثبت یکنوا ترتیب مقدارها را حفظ می‌کند و آستانه‌ها با آن تبدیل می‌شوند؛ اما فاصلهٔ نسبی ویژگی‌ها در kNN می‌تواند عوض شود.}
\end{selftest}
\begin{selftest}{رأی و انتخاب تصادفی ویژگی}
چهار رأی مثبت و سه منفی چه خروجی‌ای می‌دهند؟ آیا سهم رأی، احتمال کالیبره است؟ آیا زیرمجموعهٔ ویژگی‌های نامزد در کل یک درخت ثابت است؟
\answer{خروجی مثبت است؛ سهم $4/7$ تضمین کالیبراسیون ندارد. در جنگل متعارف، ویژگی‌های نامزد در هر گره دوباره نمونه‌گیری می‌شوند، نه فقط یک‌بار برای کل درخت.}
\end{selftest}
\begin{selftest}{محاسبهٔ معیارها از چهار شمارش}
شمارش‌های $TP=8,FN=2,FP=4,TN=86$ را در نظر بگیرید.
\answer{\begin{align*}
\operatorname{Accuracy}&=0.94,&\operatorname{Precision}&=\frac23,\\
\operatorname{Recall}&=0.8,&F_1&=\frac8{11}\approx0.7273,\\
\operatorname{FPR}&=\frac4{90}\approx0.04444.
\end{align*}
آیا مدل با صحت ۹۹ درصد که هیچ مثبت نمی‌گیرد حتماً بهتر است؟ به معیار و هزینه بستگی دارد؛ برای کشف مثبت قابل دفاع نیست.}
\end{selftest}
\begin{selftest}{نشت و انتخابِ دوباره را تشخیص دهید}
چرا محاسبهٔ مقیاس پیش از CV نشت است؟ چرا انتخاب بهترین دور بیرونی مجاز نیست؟ در طرح ۴ بیرونی، ۳ داخلی و ۲۰ تنظیم چند برازش لازم است؟
\answer{آمار بخش کنارگذاشته‌شده وارد آموزش شده است؛ انتخاب بهترین دور بیرونی نیز انتخاب بر اساس آزمون است. با بازبرازش هر دور، تعداد $4(3\times20+1)=244$ است.}
\end{selftest}
\end{selftests}

% !TeX root = ../main.tex

\end{document}
